\documentclass[10pt,a4paper]{amsart}
\usepackage{amsmath,amssymb,amsthm,graphicx}
% Class-compatibility preamble for a reader excerpt of
%   Paul Balmer, "Homological support of big objects in tensor-triangulated categories",
%   Journal de l'Ecole polytechnique -- Mathematiques 7 (2020), DOI 10.5802/jep.135.
%
% This file is the source directory's inherited layout adapter, unchanged except for one
% mechanical line: \smaller is declared with \providecommand instead of \newcommand, so that
% the excerpt compiles whether or not the wrapper class already defines \smaller as a
% font-size command. No mathematical macro, symbol, statement, proof step or reference was
% added, changed or removed by this file.
\usepackage{amsmath,amssymb,amsthm,xfrac,xspace,hyperref,graphicx,comment}
\excludecomment{DefTralics}
\providecommand{\C}{}
\providecommand{\bibliofont}{\small}
\newcommand{\ie}{i.e.,\xspace}
\newcommand{\eg}{e.g.,\xspace}
\newcommand{\etc}{etc.\xspace}
\newcommand{\cf}{cf.\xspace}
\newcommand{\resp}{resp.\xspace}
\providecommand{\smaller}{\small}
\newcommand{\Subsection}{\subsection}
\newcommand{\Subsubsection}{\subsubsection}
\newcommand{\skpt}{\smallskip}
\emergencystretch=3em
\providecommand{\loccit}{\textit{loc.\ cit.}}
\newenvironment{enumeratea}
{\bgroup\def\theenumi{\alph{enumi}}\begin{enumerate}}
{\end{enumerate}\egroup}
\newenvironment{enumeratei}
{\bgroup\def\theenumi{\roman{enumi}}\begin{enumerate}}
{\end{enumerate}\egroup}

\def\hto{\mathrel{\lhook\joinrel\to}}
\def\mto{\mathchoice{\longmapsto}{\mapsto}{\mapsto}{\mapsto}}
\def\To#1{\mathchoice{\xrightarrow{\textstyle\kern4pt#1\kern3pt}}{\stackrel{#1}{\longrightarrow}}{}{}}

\newcommand{\onto}{\mathchoice{\to\hspace*{-5mm}\to}{\to\hspace*{-2.85mm}\to}{\to\hspace*{-3mm}\to}{\to\hspace*{-3mm}\to}}

\newcommand{\into}{\mathchoice{\xymatrix@=.5cm{\ar@{>->}[r]&}}
{\xymatrix@=.3cm{\ar@{>->}[r]&}\hspace*{-1mm}}
{\hspace*{-1mm}\xymatrix@=.25cm{\ar@{>->}[r]&}\hspace*{-1.5mm}}
{\xymatrix@=.1cm{\ar@{>->}[r]&}\hspace*{-1mm}}
}
\usepackage[mathscr]{eucal}
\usepackage{bbm}
\usepackage{cleveref}
\crefdefaultlabelformat{#2\textup{#1}#3}
\usepackage[all,cmtip]{xy}\let\labelstyle\textstyle
\newdir{ >}{{}*!/-10pt/\dir{>}}
\newdir^{ (}{{}*!/-5pt/\dir^{(}}
\newdir_{ (}{{}*!/-5pt/\dir_{(}}

\numberwithin{equation}{section}

\swapnumbers

\newtheorem{Thm}[equation]{Theorem}
\newtheorem*{Thm*}{Theorem}
\newtheorem{Prop}[equation]{Proposition}
\newtheorem{Lem}[equation]{Lemma}
\newtheorem{Cor}[equation]{Corollary}

\theoremstyle{definition}
\newtheorem{Def}[equation]{Definition}
\newtheorem{Exa}[equation]{Example}
\newtheorem{Exas}[equation]{Examples}
\newtheorem{Cons}[equation]{Construction}
\newtheorem{Hyp}[equation]{Hypothesis}
\newtheorem{Rem}[equation]{Remark}
\newtheorem{Rec}[equation]{Recollection}

\DeclareMathOperator{\coker}{coker}
\DeclareMathOperator{\colim}{colim}
\DeclareMathOperator{\cone}{cone}
\DeclareMathOperator{\End}{End}
\DeclareMathOperator{\Hom}{Hom}
\DeclareMathOperator{\id}{id}
\DeclareMathOperator{\Id}{Id}
\DeclareMathOperator{\Ind}{Ind}
\DeclareMathOperator{\Ker}{Ker}
\DeclareMathOperator{\opname}{op}
\DeclareMathOperator{\supp}{supp}
\DeclareMathOperator{\Spc}{Spc}
\DeclareMathOperator{\Spec}{Spec}

\newcommand{\bbZ}{\mathbb{Z}}
\newcommand{\cat}[1]{\mathscr{#1}}
\newcommand{\Endcat}[1]{\End_{\cat #1}}
\newcommand{\Homcat}[1]{\Hom_{\cat #1}}
\newcommand{\inv}{^{-1}}
\newcommand{\isoto}{\overset{\sim}{\longrightarrow}}
\newcommand{\isotoo}{\overset{\sim}{\longrightarrow}}
\newcommand{\op}{^{\opname}}
\let\xto\To
\newcommand{\oto}[1]{\overset{#1} \longrightarrow}
\newcommand{\otoo}[1]{\overset{#1}{\longrightarrow}}
\newcommand{\Paul}[1]{{\color{ForestGreen}#1}}
\newcommand{\SET}[2]{\big\{#1\,\big|\,#2\big\}}
\newcommand{\too}{\mathop{\longrightarrow}\limits}
\newcommand{\unit}{\mathbbm{1}}

\DeclareMathOperator{\Ab}{Ab}
\DeclareMathOperator{\Add}{Add}
\DeclareMathOperator{\coev}{coev}
\DeclareMathOperator{\Der}{D}
\DeclareMathOperator{\Img}{Im}
\DeclareMathOperator{\Loc}{Loc}
\DeclareMathOperator{\modname}{mod}
\DeclareMathOperator{\Mod}{Mod}
\DeclareMathOperator{\SH}{SH}
\newcommand{\smallh}{\mathrm{h}}
\DeclareMathOperator{\Stab}{Stab}
\DeclareMathOperator{\Supp}{Supp}

\newcommand{\ababs}{\textit{ab absurdo}}
\newcommand{\ABp}{\bar{\cA}_{\cB'}}
\newcommand{\AB}{\bar{\cA}_{\cB}}
\newcommand{\AC}{\bar{\cA}_{\cC}}
\newcommand{\adj}{\dashv}
\newcommand{\bat}[1]{\bar{\cat{#1}}}
\newcommand{\bEB}{\bar E_{\cB}}
\newcommand{\bEC}{\bar E_{\cC}}
\newcommand{\boneda}{\bar\yoneda}
\newcommand{\bunit}{\bar\unit}
\newcommand{\cA}{\cat{A}}
\newcommand{\cB}{\cat{B}}
\newcommand{\cC}{\cat{C}}
\newcommand{\cD}{\cat{D}}
\newcommand{\cP}{\cat{P}}
\newcommand{\cS}{\cat{S}}
\newcommand{\cT}{\cat{T}}
\newcommand{\EBp}{E_{\cB'}}
\newcommand{\EB}{E_{\cB}}
\newcommand{\EC}{E_{\cC}}
\newcommand{\fp}{^{\textup{fp}}}
\newcommand{\hBp}{\boneda_{\cB'}}
\newcommand{\hB}{\boneda_{\cB}}
\newcommand{\hC}{\boneda_{\cC}}
\newcommand{\hEB}{\hat E_{\cB}}
\newcommand{\hEC}{\hat E_{\cC}}
\let\hook\hto
\newcommand{\hunit}{\hat\unit}
\newcommand{\ideal}[1]{\langle #1\rangle}
\newcommand{\modd}{\modname\text{\!-}}
\newcommand{\Modd}{\Mod\text{\!-}}
\newcommand{\MS}{\Modd\cat S^c}
\newcommand{\mS}{\modd\cat S^c}
\newcommand{\MT}{\Modd\cat T^c}
\newcommand{\mT}{\modd\cat T^c}
\newcommand{\potimes}[1]{^{\otimes #1}}
\newcommand{\qquadtext}[1]{\qquad\text{#1}\qquad}
\newcommand{\restr}[1]{_{|_{\scriptstyle #1}}}
\newcommand{\sbull}{{\scriptscriptstyle\bullet}}
\newcommand{\SpchS}{\Spch(\cat{S}^c)}
\newcommand{\SpchT}{\Spch(\cat{T}^c)}
\newcommand{\Spch}{\Spc^{\smallh}}
\newcommand{\SpcS}{\Spc(\cat{S}^c)}
\newcommand{\SpcT}{\Spc(\cat{T}^c)}
\newcommand{\tensids}{$\otimes$\nobreakdash-ideals}
\newcommand{\tensid}{$\otimes$\nobreakdash-ideal}
\newcommand{\tensnil}{$\otimes$\nobreakdash-nilpotent}
\newcommand{\vcorrect}[1]{{\vphantom{\vbox to #1em{}}}}
\newcommand{\yoneda}{\mathbbm{h}}
\crefname{originallemma}{original lemma}{original lemmas}
\Crefname{originallemma}{Original lemma}{Original lemmas}
\makeatletter
\newlabel{Cons:E_C}{{6}{source}{}{sourcepointer-1}{}}
\newlabel{Cons:E_C@cref}{{[originallemma][0][]6}{[1][0][]source}{}{}{}}
\makeatother
\crefname{originallemma}{original lemma}{original lemmas}
\Crefname{originallemma}{Original lemma}{Original lemmas}
\makeatletter
\newlabel{Prop:A/C}{{4}{source}{}{sourcepointer-2}{}}
\newlabel{Prop:A/C@cref}{{[originallemma][0][]4}{[1][0][]source}{}{}{}}
\makeatother
\crefname{originallemma}{original lemma}{original lemmas}
\Crefname{originallemma}{Original lemma}{Original lemmas}
\makeatletter
\newlabel{Prop:Finv}{{3}{source}{}{sourcepointer-3}{}}
\newlabel{Prop:Finv@cref}{{[originallemma][0][]3}{[1][0][]source}{}{}{}}
\makeatother
\crefname{originallemma}{original lemma}{original lemmas}
\Crefname{originallemma}{Original lemma}{Original lemmas}
\makeatletter
\newlabel{Prop:ring-in-AB}{{5}{source}{}{sourcepointer-4}{}}
\newlabel{Prop:ring-in-AB@cref}{{[originallemma][0][]5}{[1][0][]source}{}{}{}}
\makeatother
\crefname{originallemma}{original lemma}{original lemmas}
\Crefname{originallemma}{Original lemma}{Original lemmas}
\makeatletter
\newlabel{Rec:T-S}{{1}{source}{}{sourcepointer-5}{}}
\newlabel{Rec:T-S@cref}{{[originallemma][0][]1}{[1][0][]source}{}{}{}}
\makeatother
\crefname{originalequation}{original equation}{original equations}
\Crefname{originalequation}{Original equation}{Original equations}
\makeatletter
\newlabel{eq:hatF-hatU}{{1}{source}{}{sourcepointer-6}{}}
\newlabel{eq:hatF-hatU@cref}{{[originalequation][0][]1}{[1][0][]source}{}{}{}}
\makeatother
\crefname{originallemma}{original lemma}{original lemmas}
\Crefname{originallemma}{Original lemma}{Original lemmas}
\makeatletter
\newlabel{it:A/C-d}{{4}{source}{}{sourcepointer-7}{}}
\newlabel{it:A/C-d@cref}{{[originallemma][0][]4}{[1][0][]source}{}{}{}}
\makeatother
\title{The Serre tensor ideal cut out by a homological functor, and the localisation diagram it produces}
\author{Paul Balmer}
\date{doi:10.5802/jep.135; CC BY4.0}
\begin{document}
\maketitle
\noindent\textit{Editorial note.} All delivered mathematical text is the author's own, copied from the pinned publisher source sources/jep-2020-135-homological-support/full-original-native.tex; no statement or proof step was rewritten, reordered or removed.\par\medskip
\noindent\textit{Editorial note.} Preamble provenance: the retained preamble is the source's own layout adapter plus the author's macro block at lines 6-143 of sources/jep-2020-135-homological-support/full-original-native.tex (everything between the publisher class metadata and \textbackslash{}begin\{document\}, comments excluded). Nothing was added, no macro was invented, and the publisher document class is deliberately not reproduced; the excerpt serves only as preamble, never as a span coordinate.\par\medskip
\noindent\textit{Editorial note.} External references are navigation pointers only; every deliberately kept step of the author's argument stays inside the counted spans. The displayed numbers of those pointers are the author's own numbering under the author's \textbackslash{}newtheorem declarations (the source has no subsections, so numbering is section.index), computed mechanically by the lane's scanner.\par\medskip
\noindent\textit{Editorial note.} The counted statement includes the Xy-pic display of the four-square diagram, and the proof's first sentence notes that the left-hand square is a repetition for the reader's convenience, which is kept verbatim.\par\medskip
\noindent\textit{Editorial note.} Every \textbackslash{}Cref outside the counted range is recorded as a navigation stub whose displayed number is the author's own numbering computed from the author's \textbackslash{}newtheorem declarations; the load-bearing steps of the proof stay inside the counted spans.\par\medskip
\noindent\textit{Editorial note.} Assembly history: the first attempt (r001) was returned with 'invalid external label' because the counted proof refers to the author's labels Prop:A/C and it:A/C-d, which contain a slash, and the assembler's external-label whitelist did not admit that character. The Chief fixed the tool on 2026-10-12 (the whitelist now admits '/', while characters that could break out of the label such as braces, backslash, hash, percent and tilde stay refused), so the card is reinstated unchanged and re-assembled as r002; no statement, span or stub changed. The withdrawn job directory was moved back into jobs/author-material-new/ on the local mirror before this run.\par\medskip
\noindent\textit{Allowed standard prerequisites (editorial).}\begin{itemize}
\item The standing hypothesis on a big tensor-triangulated category and its tensor-triangulated subcategory of compact objects. Same source, Hypothesis labelled Hyp:big-tt (source lines 180-183); invoked in the counted statement. Fixes the ambient category mT and the class of objects the construction applies to.
\item The homological functor criterion producing a maximal Serre tensor ideal. Same source, Proposition labelled Prop:Finv; invoked inside the counted proof at line 799. Supplies the maximal Serre tensor ideal cB = Ker(hat H) cap A\textasciicircum{}fp from the homological functor H = Q\_cC o hat F o yoneda\_T.
\item Localisation of a tensor-triangulated category at a Serre tensor ideal. Same source, the notation AB = mT/Loc(cB) set up before the statement; used at lines 797-806. Produces the quotient category and the localisation functor Q\_cB through which the homological functor factors.
\end{itemize}\medskip
\ifdefined\hypertarget\hypertarget{sourcepointer-1}{}\fi
\noindent\textit{Source reference.} Original lemma 6: sources/jep-2020-135-homological-support/full-original-native.tex, source lines 784--784. Navigation pointer for \textbackslash{}ref\{Cons:E\_C\} used inside the counted span (line 784); the argument that consumes it is retained in the span.\par\medskip
\ifdefined\hypertarget\hypertarget{sourcepointer-2}{}\fi
\noindent\textit{Source reference.} Original lemma 4: sources/jep-2020-135-homological-support/full-original-native.tex, source lines 814--814. Navigation pointer for \textbackslash{}ref\{Prop:A/C\} used inside the counted span (line 814); the argument that consumes it is retained in the span.\par\medskip
\ifdefined\hypertarget\hypertarget{sourcepointer-3}{}\fi
\noindent\textit{Source reference.} Original lemma 3: sources/jep-2020-135-homological-support/full-original-native.tex, source lines 799--799. Navigation pointer for \textbackslash{}ref\{Prop:Finv\} used inside the counted span (line 799); the argument that consumes it is retained in the span.\par\medskip
\ifdefined\hypertarget\hypertarget{sourcepointer-4}{}\fi
\noindent\textit{Source reference.} Original lemma 5: sources/jep-2020-135-homological-support/full-original-native.tex, source lines 814--814. Navigation pointer for \textbackslash{}ref\{Prop:ring-in-AB\} used inside the counted span (line 814); the argument that consumes it is retained in the span.\par\medskip
\ifdefined\hypertarget\hypertarget{sourcepointer-5}{}\fi
\noindent\textit{Source reference.} Original lemma 1: sources/jep-2020-135-homological-support/full-original-native.tex, source lines 767--767. Navigation pointer for \textbackslash{}ref\{Rec:T-S\} used inside the counted span (line 767); the argument that consumes it is retained in the span.\par\medskip
\ifdefined\hypertarget\hypertarget{sourcepointer-6}{}\fi
\noindent\textit{Source reference.} Original equation 1: sources/jep-2020-135-homological-support/full-original-native.tex, source lines 788--788. Navigation pointer for \textbackslash{}ref\{eq:hatF-hatU\} used inside the counted span (line 788); the argument that consumes it is retained in the span.\par\medskip
\ifdefined\hypertarget\hypertarget{sourcepointer-7}{}\fi
\noindent\textit{Source reference.} Original lemma 4: sources/jep-2020-135-homological-support/full-original-native.tex, source lines 814--814. Navigation pointer for \textbackslash{}ref\{it:A/C-d\} used inside the counted span (line 814); the argument that consumes it is retained in the span.\par\medskip
\begin{Lem}
\label{Lem:Finv}
With notation as in~\Cref{Rec:T-S}, let $\cC\in\SpchS$.
Then $\cB:=\SET{M\in\mT}{\hat F(M)\in\cC}$ is a maximal Serre \tensid\ of~$\mT$. Furthermore, we have a diagram
\begin{equation}
\label{eq:F-res-fields}
\vcenter{\xymatrix@R=1.5em{
\cT \ar@<-.3em>[dd]_-{F} \ar[r]^-{\yoneda_\cT}
& \cA=\MT {}_{\vcorrect{.6}} \ar@<-.3em>[dd]_-{\hat F} \ar@{->>}@<-.2em>[rd]_-{Q_{\cB}}
\\
&& \AB=\MT/\Loc(\cB) {}^{\vcorrect{.6}} \ar@<-.3em>[dd]_-{\bar F} \ar@{_{ (}->}@<-.2em>[lu]_-{R_{\cB}}
\\
\cS \ar@<-.3em>[uu]_-{U} \ar[r]^-{\yoneda_\cS}
& \MS {}_{\vcorrect{.6}} \ar@<-.3em>[uu]_-{\hat U} \ar@{->>}@<-.2em>[rd]_-{Q_{\cC}}
\\
&& \MS/\Loc(\cC) {}^{\vcorrect{.6}} \ar@<-.3em>[uu]_-{\bar U} \ar@{_{ (}->}@<-.2em>[lu]_-{R_{\cC}}
}}
\end{equation}
with $\yoneda_\cS\,F\cong\hat F\,\yoneda_\cT$ and $\yoneda_\cT\,U\cong\hat U\,\yoneda_\cS$ and $Q_{\cC}\,\hat F\cong\bar F\,Q_{\cB}$ and $\hat U\,R_{\cC}\cong R_{\cB}\,\bar U$.
Finally, the pure-injective $\EB$ is a direct summand of~$U(\EC)$ in~$\cT$ (see \Cref{Cons:E_C}).
\end{Lem}

\begin{proof}
The left-hand square of~\eqref{eq:F-res-fields} is~\eqref{eq:hatF-hatU}, repeated for cognitive help. Since $\hat F$ is exact, $\hat U$ preserves injectives hence $U(\EC)$ is pure-injective. Since $U$ is lax-monoidal, $U(\EC)$ is a weak ring, with unit~$\unit\to U(\EC)$ adjoint to $\eta_{\cC}\colon \unit_\cS=F(\unit_\cT)\to \EC$. In particular, $U(\EC)\neq 0$. The internal-hom version of the adjunction reads
\begin{equation}
\label{eq:aux-internal-adj}
[M,\hat U(\hEC)]\cong \hat U[\hat F(M),\hEC]
\end{equation}
for every~$M\in\cA$, as can be checked by testing under $\Homcat{A}(\hat x,-)$ for $x\in\cT^c$.

Consider now the functor $H:=Q_{\cC}\circ \hat F\circ\yoneda_\cT$
\[
H\colon \cT\to \MT\to \MS\onto\MS/\Loc(\cC)=:\cD.
\]
This is a homological functor to which we can apply \Cref{Prop:Finv}, with $\hat H$ being necessarily~$Q_{\cC}\circ\hat F$. It tells us that $\cB=\Ker(\hat H)\cap\cA\fp$ is indeed a maximal Serre \tensid\ of~$\cA\fp$. We can then factor~$\hat H$ via~$Q_\cB$, yielding a unique tensor-exact coproduct-preserving functor $\bar F\colon \MT/\Loc(\cB)\to \MS/\Loc(\cC)$ making the following diagram commute
\[
\xymatrix@=1.5cm{
\MT \ar[d]_-{\hat F} \ar@{->>}[r]^-{Q_{\cB}} \ar[rd]_-{\hat H}
& \MT/\Loc(\cB)=\AB \ar@{..>}[d]^-{\bar F}_-{\exists\,!}
\\
\MS \ar@{->>}[r]_-{Q_{\cC}}
& \MS/\Loc(\cC) \ar@{..>}@/_3em/[u]_-{\bar U}
}
\]
By general Grothendieck-category theory, $\bar F$ admits a right adjoint~$\bar U$. This gives us~\eqref{eq:F-res-fields}. Note that $R_{\cB}\bar U$ is the right adjoint to~$\bar F\,Q_{\cB}\cong Q_{\cC}\,\hat F$, whose right adjoint is also~$\hat U\,R_{\cC}$. Hence we have
\begin{equation}
\label{eq:aux-RU}
R_{\cB}\circ \bar U \cong \hat U\circ R_{\cC}.
\end{equation}
Consider $\bunit\into \bEC$ the injective hull of the unit in~$\MS/\Loc(\cC)$, as in the statement. The pure-injective $\EC\in\cS$ corresponding to~$\cC$ is then characterized by~$\hEC=R_{\cC}(\bEC)$. Therefore, by~\eqref{eq:aux-RU}, we have $R_{\cB}(\bar U(\bEC))\cong \hat U(\hEC)=\yoneda_{\cT}(U(\EC))$. Applying~$Q_{\cB}$ to this relation, we see that the injective $\bar U(\bEC)$ is the image in~$\AB$ of~$U(\EC)$. By \Cref{Prop:ring-in-AB} the unit of this non-trivial weak ring $\bar U(\bEC)$ is a monomorphism $\bunit\into \bar U(\bEC)$ in~$\AB$ and $\bar U(\bEC)$ is injective. Therefore the injective hull $\bEB$ of~$\bunit$ is a direct summand of~$\bar U(\bEC)$. As usual, this decomposition holds in~$\cT$ already, by \Cref{Prop:A/C}\,\eqref{it:A/C-d}.
\end{proof}

\begin{thebibliography}{99}
\bibitem{Balmer05a} P. Balmer - “The spectrum of prime ideals in tensor triangulated categories” , J. reine angew. Math. 588 (2005), p. 149-168
\bibitem{Balmer18a} P. Balmer - “On the surjectivity of the map of spectra associated to a tensor-triangulated functor” , Bull. London Math. Soc. 50 (2018) no. 3, p. 487-495
\bibitem{Balmer19} P. Balmer - “A guide to tensor-triangular classification” , in Handbook of homotopy theory (H. Miller, ed.), Chapman and Hall/CRC, 2019, Available on the author’s web page
\bibitem{Balmer20a} P. Balmer - “Nilpotence theorems via homological residue fields” , Tunis. J. Math. 2 (2020) no. 2, p. 359-378
\bibitem{BalmerCameron20pp} P. Balmer \& J. Cameron - “Computing homological residue fields in algebra and topology” , 2020
\bibitem{BalmerDellAmbrogioSanders16} P. Balmer, I. Dell’Ambrogio \& B. Sanders - “Grothendieck-Neeman duality and the Wirthmüller isomorphism” , Compositio Math. 152 (2016) no. 8, p. 1740-1776
\bibitem{BalmerFavi11} P. Balmer \& G. Favi - “Generalized tensor idempotents and the telescope conjecture” , Proc. London Math. Soc. (3) 102 (2011) no. 6, p. 1161-1185
\bibitem{BensonIyengarKrause08} D. J. Benson, S. B. Iyengar \& H. Krause - “Local cohomology and support for triangulated categories” , Ann. Sci. École Norm. Sup. (4) 41 (2008) no. 4, p. 573-619
\bibitem{BensonIyengarKrause11a} D. J. Benson, S. B. Iyengar \& H. Krause - “Stratifying modular representations of finite groups” , Ann. of Math. (2) 174 (2011) no. 3, p. 1643-1684
\bibitem{BensonIyengarKrause11c} D. J. Benson, S. B. Iyengar \& H. Krause - “Stratifying triangulated categories” , J. Topology 4 (2011) no. 3, p. 641-666
\bibitem{BensonIyengarKrause12b} D. J. Benson, S. B. Iyengar \& H. Krause - “Colocalizing subcategories and cosupport” , J. reine angew. Math. 673 (2012), p. 161-207
\bibitem{BensonIyengarKrause12a} D. J. Benson, S. B. Iyengar \& H. Krause - Representations of finite groups: local cohomology and support , Oberwolfach Seminars , vol. 43 , Birkhäuser/Springer, Basel, 2012
\bibitem{BensonIyengarKrause13} D. J. Benson, S. B. Iyengar \& H. Krause - “Module categories for group algebras over commutative rings” , J. K-Theory 11 (2013) no. 2, p. 297-329, With an appendix by Greg Stevenson
\bibitem{BalmerKrauseStevenson19} P. Balmer, H. Krause \& G. Stevenson - “Tensor-triangular fields: ruminations” , Selecta Math. (N.S.) 25 (2019) no. 1, article ID 13, 36 pages
\bibitem{BalmerKrauseStevenson20} P. Balmer, H. Krause \& G. Stevenson - “The frame of smashing tensor-ideals” , Math. Proc. Cambridge Philos. Soc. 168 (2020) no. 2, p. 323-343
\bibitem{DwyerPalmieri08} W. G. Dwyer \& J. H. Palmieri - “The Bousfield lattice for truncated polynomial algebras” , Homology Homotopy Appl. 10 (2008) no. 1, p. 413-436
\bibitem{HoveyPalmieriStrickland97} M. Hovey, J. H. Palmieri \& N. P. Strickland - Axiomatic stable homotopy theory , Mem. Amer. Math. Soc. , vol. 128, no. 610 , American Mathematical Society, Providence, RI, 1997
\bibitem{HoveyStrickland99} M. Hovey \& N. P. Strickland - Morava $K$-theories and localisation , Mem. Amer. Math. Soc. , vol. 139, no. 666 , American Mathematical Society, Providence, RI, 1999
\bibitem{Krause00} H. Krause - “Smashing subcategories and the telescope conjecture—an algebraic approach” , Invent. Math. 139 (2000) no. 1, p. 99-133
\bibitem{LurieHA} J. Lurie - “Higher algebra” (2017), Online at http://www.math.ias.edu/\textasciitilde{}lurie/
\bibitem{Neeman96} A. Neeman - “The Grothendieck duality theorem via Bousfield’s techniques and Brown representability” , J. Amer. Math. Soc. 9 (1996) no. 1, p. 205-236
\bibitem{Neeman00} A. Neeman - “Oddball Bousfield classes” , Topology 39 (2000) no. 5, p. 931-935
\bibitem{Neeman01} A. Neeman - Triangulated categories , Annals of Math. Studies , vol. 148 , Princeton University Press, Princeton, NJ, 2001
\bibitem{Stevenson13} G. Stevenson - “Support theory via actions of tensor triangulated categories” , J. reine angew. Math. 681 (2013), p. 219-254
\end{thebibliography}
\end{document}
